\documentclass[aspectratio=169]{beamer}
\usepackage[utf8]{inputenc}\usepackage{amsmath}\usepackage{tcolorbox}
\definecolor{UNIBBlue}{RGB}{0, 51, 102}\setbeamercolor{frametitle}{bg=UNIBBlue,fg=white}
\title{Optimisasi untuk Teknik Elektro}\subtitle{Minggu 12: Particle Swarm Optimization (PSO)}
\date{Semester Ganjil 2026}\begin{document}\begin{frame}\titlepage\end{frame}
\begin{frame}{PSO: Konsep Dasar}
Terinspirasi dari pergerakan kawanan burung atau ikan.
Setiap kandidat solusi adalah "partikel" yang memiliki kecepatan dan posisi di ruang solusi.
Update Kecepatan dipengaruhi oleh:
\begin{itemize}
\item Inersia (mempertahankan arah saat ini).
\item \textit{Personal Best} / pbest (pengalaman memori terbaik individu).
\item \textit{Global Best} / gbest (pengalaman memori terbaik kelompok).
\end{itemize}
\end{frame}\end{document}